\scanpage{7}
Si $\beta$ inversible on conclut $a=\lambda\pi$, donc l'équation qui est dessus \inserted{(\unsure{*|8})} devient
\[
\begin{aligned}
& f(a)\pi'+\lambda\underbrace{\psi(\pi)}_{-\sigma\pi'}=0\qquad\text{i.e.}\quad \sigma\lambda=\lambda^2\underbrace{f(\pi)}_{-\widetilde\sigma\widetilde\rho}\\
&\deleted{\beta\lambda\pi=\sigma'\lambda^2 f(\pi)\pi}\\
\text{i.e.}\quad &\beta\lambda=\sigma'\lambda^2\underbrace{f(\pi)}_{-\widetilde\sigma\widetilde\rho}\quad\deleted{\text{i.e.}}\qquad\text{i.e.}\quad\lambda(\sigma+\lambda\widetilde\sigma\widetilde\rho)=0\\
\text{i.e.}\quad &\lambda(\beta+\sigma'\widetilde\sigma\widetilde\rho\lambda)=0,\quad\text{en plus de }\rho\lambda=0
\end{aligned}
\]
Si on est dans un anneau intègre, celà implique \uline{soit} $\lambda=0$ i.e. $u=\mathrm{id}$, soit $\beta+\sigma'\widetilde\sigma\widetilde\rho\lambda=0$ ce qui implique (comme $\beta$ inv. par hyp) $\sigma'\widetilde\sigma\widetilde\rho$ inversible (et $\lambda=-\frac{\deleted{\unsure{\sigma'\widetilde\sigma\widetilde\rho}}\beta}{\sigma'\widetilde\sigma\widetilde\rho}$) absurde si $\rho=\bar\rho$ ($n$ impair) qui implique $\lambda=0$ (puisque $\rho\lambda=0$), sinon $\rho=2\widetilde\rho\deleted{\unsure{n}}$, si $2$ inv. \deleted{même \unsure{concl.} (\unsure{même} résultat)} même résultat. Mais si car. $2$ et $n$ pair, $\deleted{\beta\lambda=0}$ $\rho=0$ donc $\rho\lambda=0$ par des \unsure{ventres bas|ventribus|en conclure}, il reste la condition (\unsure{néc.} et suff$^{\mathrm{te}}$) $\sigma+\lambda\widetilde\sigma\widetilde\rho=0\deleted{\unsure{+}}$. Or $\sigma\sigma'=\beta+\gamma\rho=\beta$ \deleted{\unsure{donc}} inv, donc \inserted{$\sigma,\sigma'$ et} $\widetilde\sigma=\deleted{\unsure{\widetilde{\phantom\sigma}}}\sigma$ \unsure{est} inv, la condition devient $1+\lambda\widetilde\rho=0$ i.e. $\widetilde\rho$ inv. et $\lambda=-\frac1{\widetilde\rho}$. Ainsi

\uline{Prop.} Sur anneau $k$ intègre, si $\beta$ ou $\rho$ $\ne0$, alors le seul $u\in\operatorname{Aut}(E,f,\pi,\pi')$ induisant l'identité dans $V$ est l'identité, \uline{sauf} dans le cas où $\deleted{\rho=0}$, $\beta\ne0$, $n$ pair, $\deleted{\rho=2\rho'}$ car. $2$ $(\rho=2\rho'=0)$, $\rho'$ inv. auquel cas le seul autre $u$ est $\mathrm{id}\unsure{+|-}\frac1{\widetilde\rho}\pi'\otimes\pi=\mathrm{id}+\frac1{\widetilde\rho}\unsure{\langle\pi',\pi'\rangle|\langle\pi,\pi'\rangle}$ i.e. l'antipodisme

% archive sheet number: 6

